% 本稿只演示引用与研究计划的衔接，未执行TeX编译，不是正式投稿稿件。
% 三个引用键在同目录references.bib中存在；BERT和Longformer仅核对元数据与摘要。
\documentclass{article}
\usepackage[T1]{fontenc}
\title{Attention Connectivity: A Reading and Study Proposal}
\author{Workflow Example}
\date{}
\begin{document}
\maketitle

\section{Related Work and Evidence Boundary}
The Transformer organizes sequence transduction around attention~\cite{vaswani2017attention}.
BERT studies bidirectional language representation pre-training~\cite{devlin2018bert},
whereas Longformer addresses long documents using local and global attention~\cite{beltagy2020longformer}.
These are different research settings; their reported task scores do not form a common ranking.
Our reading of the latter two sources is limited to their abstracts.
Their full methodological and evaluation details remain to be checked.

\section{Proposed Study}
% 使用将来时和提案措辞，不把未执行的实验写成已经观察到的结论。
We propose a controlled key--value retrieval study with dense, windowed,
and windowed-plus-global connectivity under a fixed small-model budget.
We will report accuracy by sequence length and retrieval distance.
Runtime and memory measurements will be interpreted according to the actual
implementation; a dense masked kernel will not be described as sparse acceleration.

\section{Current Status}
No implementation, training, timing measurements, or reproduction has been completed.
Consequently, we report no numerical results or performance improvements.
The experiment specification and empty result schema are maintained in the accompanying Chinese note.

\bibliographystyle{plain}
\bibliography{references}
\end{document}
